\documentclass[11pt]{article}
\usepackage[margin=1.1in]{geometry}
\usepackage{amsmath,amssymb,amsthm}
\usepackage{booktabs}
\usepackage{graphicx}
\usepackage[colorlinks=true,linkcolor=blue]{hyperref}

\newcommand{\uu}{\boldsymbol{u}}
\newcommand{\uh}{\uu_h}
\newcommand{\om}{\boldsymbol{\omega}}
\newcommand{\jump}[1]{[\![#1]\!]}
\newcommand{\avg}[1]{\{\!\!\{#1\}\!\!\}}
\newcommand{\dd}{\mathrm{d}}
\newcommand{\Rd}{R_d}
\newcommand{\cv}{c_v}
\newcommand{\cp}{c_p}
\newcommand{\Ja}[1]{(J\boldsymbol{a}^{#1})}
\newcommand{\Fs}{F^\sharp}

\newtheorem{prop}{Proposition}
\newtheorem{rem}{Remark}

\title{Mathematics of the Vector-Invariant DG--FD Dry Dynamical Core\\
{\large \texttt{experiments/dg\_dycore} --- kinetic-energy-preserving face
terms, entropy-dissipative interfaces, well-balanced Exner pressure
gradients, and terrain verification}}
\author{ClimaAtmos DG sandbox}
\date{August 2026}

\begin{document}
\maketitle

\begin{abstract}
This note collects the continuous equations, the discrete operators, and
the main structure-preservation results implemented and verified in the
\texttt{dg\_dycore} sandbox: a horizontally discontinuous-Galerkin (DG),
vertically staggered finite-difference (FD) dry compressible core in
vector-invariant (VI) form, sharing its metric machinery with the
continuous-Galerkin ClimaAtmos dycore. Every discrete operator is written
out at the nodal level, in the constructive style of flux-differencing DG
expositions. The central results are (i) the unique two-point face-flux
set for which the discrete horizontal advective kinetic-energy budget
closes exactly (KEP), on flat and terrain-warped grids alike; (ii) a
single-slot entropy-dissipative interface modification that retains exact
KEP; (iii) the reference-subtracted Exner pressure-gradient formulation
and its problem-matched reference states; and (iv) exact
discrete-hydrostatic initial states, including the effective-geopotential
generalization required over steep terrain. Verification spans
operator-level discrete budget identities (closure to $10^{-15}$
relative), a baroclinic wave over analytic double-mountain orography
(machine-precision mass and energy conservation over one day), and linear
Agnesi mountain waves reproducing $\lambda_z = 2\pi U/N$ and linear
amplitude scaling.
\end{abstract}

\section{Governing equations}
\label{sec:continuous}

The core solves the dry compressible Euler equations in a rotating frame.
The prognostic state is
\begin{equation}
Y = \big(\rho,\; \rho e,\; \uh,\; w\big),
\end{equation}
density and total energy density at cell centers, the two horizontal
covariant momentum components $\uh = (u_1, u_2)$ at cell centers, and the
vertical covariant velocity $w = u_3$ at cell faces (Lorenz-type
staggering). The continuous system is
\begin{align}
\partial_t \rho &= -\nabla\!\cdot(\rho\uu),
\label{eq:mass}\\
\partial_t (\rho e) &= -\nabla\!\cdot\big[(\rho e + p)\,\uu\big],
\label{eq:energy}\\
\partial_t \uu &= -(2\boldsymbol{\Omega} + \om)\times\uu
  \;-\; \nabla K \;-\; \frac{1}{\rho}\nabla p \;-\; \nabla\Phi,
\label{eq:vi}
\end{align}
with relative vorticity $\om = \nabla\times\uu$, kinetic energy
$K = \tfrac12|\uu|^2$, and geopotential $\Phi = gz$; the vector-invariant
identity $(\uu\cdot\nabla)\uu = \om\times\uu + \nabla K$ has been applied
in \eqref{eq:vi}. The specific total energy and the equation of state are
\begin{equation}
e = \cv (T - T_{\mathrm{tri}}) + K + \Phi,
\qquad
p = \frac{\Rd}{\cv}\Big(\rho e - \rho K - \rho\Phi
 + \rho \cv T_{\mathrm{tri}}\Big),
\label{eq:eos}
\end{equation}
an ideal gas offset by the triple-point temperature $T_{\mathrm{tri}}$,
matching the ClimaAtmos thermodynamics convention.

\subsection{Terrain-following coordinates}
\label{sec:terrain}

Orography $h(x,y)$ enters through the linear-decay warp
\begin{equation}
z(\xi^3) = \xi^3 + \Big(1 - \frac{\xi^3}{z_{\mathrm{top}}}\Big)\,h,
\end{equation}
(\texttt{Hypsography.LinearAdaption}), which defines the reference-to-physical
map $\boldsymbol{x}(\boldsymbol{\xi})$, its Jacobian determinant $J$, and
the contravariant basis $\boldsymbol{a}^d = \nabla\xi^d$. Horizontal
momentum is carried in covariant components on the warped grid. Two
metric constructions are shared verbatim with the CG ClimaAtmos core.
First, the contravariant vertical velocity at faces carries the terrain
cross-term,
\begin{equation}
\tilde u^3\big|_f
 = \underbrace{\big(C^{123}(w)\big)^3}_{\text{vertical part}}
 + \underbrace{\mathcal{I}^{\rho J}_{c\to f}
   \big[(C^{123}(\uh))^3\big]}_{\text{terrain cross-term}},
\label{eq:u3}
\end{equation}
where $\mathcal{I}^{\rho J}_{c\to f}$ is $\rho J$-weighted center-to-face
interpolation. Second, the kinetic energy uses the full metric, including
the $g^{13}, g^{23}$ cross terms that vanish on flat grids:
\begin{equation}
K = \tfrac12\Big( u_i u^i \big|_{\uh}
 + \mathcal{I}_{f\to c}\big[w_3 w^3\big]
 + 2\, u^i \,\mathcal{I}_{f\to c}[w_i] \Big).
\label{eq:fullK}
\end{equation}
Under the linear warp, $\partial \xi^h/\partial z \equiv 0$: horizontal
face normals are exactly horizontal, so the horizontal face machinery is
metric-exact and all tilt transport passes through \eqref{eq:u3}.

The kinematic lower boundary condition is imposed statically at the
bottom face only: $w_3$ is set so that $u^3 = 0$ there and is frozen by
the boundary operator. Interior $w$ is \emph{not} terrain-adapted:
interior flow legitimately crosses coordinate surfaces, and a structured
interior $w$ breaks the staggered $\nabla K$/Lamb cancellation
(Section~\ref{sec:kep}).

\section{Notation and discrete operators}
\label{sec:notation}

\subsection{Elements, nodes, and metric terms}

The horizontal mesh consists of quadrilateral elements (equiangular cubed
sphere, or a doubly periodic plane) each mapped to the reference square
$[-1,1]^2$. Within an element, fields are represented on the tensor
product of $N_q = N+1$ Gauss--Lobatto--Legendre (GLL) nodes
$\{\xi_i\}_{i=1}^{N_q}$ with quadrature weights $\{w_i\}$; the default
polynomial degree is $N = 4$. Let $D_{ik} = \ell_k'(\xi_i)$ denote the
one-dimensional differentiation matrix of the GLL Lagrange basis. Every
horizontal node carries the metric quantities
\begin{equation}
J_{ij},\qquad
\Ja{1}_{ij} = J_{ij}\,\nabla\xi^1\big|_{ij},\qquad
\Ja{2}_{ij} = J_{ij}\,\nabla\xi^2\big|_{ij},
\qquad
(WJ)_{ij} = w_i w_j J_{ij},
\end{equation}
i.e.\ the metric-weighted contravariant basis vectors and the
mass-matrix entry. The discretization is collocation-based: no
over-integration is used, and there is no continuity projection
(\texttt{dss}); elements communicate only through faces.

\subsection{Traces, jumps, averages}

On an interior face shared by elements $-$ (owner) and $+$ (neighbor),
with outward unit normal $\hat n$ of the owner, define for any nodal
quantity $q$
\begin{equation}
\jump{q} = q^- - q^+,
\qquad
\avg{q} = \tfrac12\big(q^- + q^+\big).
\end{equation}
The surface mass is obtained from the volume metric: along a constant-$\xi^1$
face, $s\!W\!J = \big| w_j \Ja{1} \big|$ and
$\hat n = \pm \Ja{1}/|\Ja{1}|$ (sign by outward orientation), and
analogously with $\Ja{2}, w_i$ on constant-$\xi^2$ faces. On extruded
grids the same construction is applied level-by-level with the product
local geometry, so $s\!W\!J$ carries the vertical measure and is
consistent with the three-dimensional $WJ$.

\subsection{The flux-differencing volume operator}
\label{sec:fd}

Let $\Fs(\boldsymbol{n}_a, \boldsymbol{n}_b;\, y_a, y_b)$ be a two-point
flux, symmetric under simultaneous exchange of its argument pairs, and
consistent: $\Fs(\boldsymbol{n},\boldsymbol{n}; y, y) =
\boldsymbol{F}(y)\cdot\boldsymbol{n}$ for the physical flux
$\boldsymbol{F}$. Note the signature: the two-point flux receives the
metric-weighted direction vectors of \emph{both} nodes --- this node-pair
metric convention is what makes the operator well-defined on curved and
warped grids. The mass-weighted volume tendency at node $(i,j)$ is the
nodal sum
\begin{equation}
\big(WJ\,\partial_t y\big)_{ij} \;+\!\!=\;
-2\,w_i w_j \sum_{k=1}^{N_q}\Big[
 D_{ik}\, \Fs\!\big(\Ja{1}_{ij}, \Ja{1}_{kj};\, y_{ij}, y_{kj}\big)
+ D_{jk}\, \Fs\!\big(\Ja{2}_{ij}, \Ja{2}_{ik};\, y_{ij}, y_{ik}\big)
\Big]
\;+\; B_{ij},
\label{eq:fdsum}
\end{equation}
followed by division by $(WJ)_{ij}$. The boundary contribution $B_{ij}$
is nonzero only on element-boundary nodes and consists of the
\emph{own-node consistent flux},
\begin{equation}
B_{ij} =
\begin{cases}
\mp\, w_j\, \Fs\big(\Ja{1}_{ij}, \Ja{1}_{ij};\, y_{ij}, y_{ij}\big),
& i = 1 \;(\,-\,)\ \text{or}\ i = N_q\ (\,+\,\text{sign flipped}),\\[2pt]
\mp\, w_i\, \Fs\big(\Ja{2}_{ij}, \Ja{2}_{ij};\, y_{ij}, y_{ij}\big),
& j = 1 \ \text{or}\ j = N_q,
\end{cases}
\end{equation}
which converts the strong-form sum into the weak-equivalent form: for
$\Fs$ consistent and constant fields, \eqref{eq:fdsum} vanishes to the
accuracy of the discrete metric identities (free-stream preservation to
truncation on curved meshes).

\subsection{Interface flux and lifting operators}

Interior faces are coupled by a simultaneous-approximation-term (SAT)
exchange. For a numerical flux $F^*(\hat n^-; y^-, y^+)$ evaluated once
per face node with the owner's normal,
\begin{equation}
\big(WJ\,\partial_t y\big)^-\; -\!\!=\; s\!W\!J \; F^*,
\qquad
\big(WJ\,\partial_t y\big)^+\; +\!\!=\; s\!W\!J\; F^*,
\label{eq:sat}
\end{equation}
an antisymmetric update that conserves $\sum WJ\, y$ identically for the
flux-form variables. Together with $B_{ij}$ of \eqref{eq:fdsum}, the pair
$(B, F^*)$ reproduces the standard weak-form surface term
$\oint \hat n \cdot (\boldsymbol{F}^* - \boldsymbol{F}^-)$.

Non-conservative (gradient, curl, penalty) terms use element-local strong
derivatives completed by \emph{lifting corrections}: for a face function
$g(\hat n;\,\cdot^-,\cdot^+)$ the correction adds, symmetrically per
side,
\begin{equation}
\big(\partial_t y\big)^\pm \;+\!\!=\;
\frac{s\!W\!J}{WJ}\; g\big(\hat n^\pm;\; \cdot\big).
\label{eq:lift}
\end{equation}
The central gradient lift $g = \tfrac12(q^+ - q^-)\,\hat n$ completes the
strong gradient of a discontinuous scalar $q$; the analogous central
curl and divergence lifts complete $\nabla\times$ and $\nabla\cdot$.
$\lambda$-scaled jump penalties use the same mechanism.

\subsection{Vertical operators and time integration}

The vertical direction is a staggered FD column: centers carry
$(\rho, \rho e, \uh)$, faces carry $w$. The operators used below are the
two-point interpolants and divergences
$\mathcal{I}_{c\to f}, \mathcal{I}_{f\to c},
\partial^{c\leftarrow f}_3, \partial^{f\leftarrow c}_3$ with the
boundary conditions stated in the source; vertical advection of scalars
uses a monotone Lin--van Leer limiter in the explicit part.

Time stepping is IMEX ARK (ARS343) in a HEVI split. The implicit residual
contains the vertical acoustic system and the Rayleigh sponge,
\begin{equation}
\partial_t \rho \supset -\partial_3(\rho w^3),\qquad
\partial_t \rho e \supset -\partial_3\big[\rho h_{\mathrm{tot}} w^3\big],\qquad
\partial_t w \supset -\frac{\partial_3 p}{\rho_f}
 - \partial_3(K+\Phi) - \beta(z)\,w,
\label{eq:implicit}
\end{equation}
solved with a fixed two-iteration Newton method (no convergence loop) and
an analytic column-block Jacobian. The sponge profile
$\beta(z) = \alpha_w \sin^2\!\big[\tfrac{\pi}{2}(z -
z_d)/(z_{\mathrm{top}}-z_d)\big]$ for $z > z_d$ (zero below; acting on
$w$ only) contributes its exact Jacobian entry $-\beta$ on the $w$--$w$
diagonal, so the ClimaAtmos damping strength $\alpha_w =
1\,\mathrm{s}^{-1}$ is unconditionally stable at any $\Delta t$.

\section{Kinetic-energy analysis and the KEP face set}
\label{sec:kep}

We analyze the semi-discrete production of the advected energy
$E = \sum_{\mathrm{nodes}} WJ\,\rho\,(K + \Phi)$ by the horizontal
advective terms: the flux-differencing mass and energy updates contracted
with $(K+\Phi)$, plus the momentum advection terms (Lamb term, $\nabla K$
and its lifting) contracted with $\rho\uu$. Call this residual
$P_{\mathrm{adv}}$; for the continuous equations $P_{\mathrm{adv}} = 0$
identically (advection neither creates nor destroys kinetic energy), and
the question is which discrete flux choices preserve the identity.

\subsection{Volume condition}

Insert \eqref{eq:fdsum} for the mass equation into
$\sum_{ij} (K+\Phi)_{ij} (WJ\,\partial_t\rho)_{ij}$ and split each
node-pair term into its symmetric and antisymmetric parts under
$(ij)\leftrightarrow(kj)$. The GLL SBP property
$w_i D_{ik} + w_k D_{ki} = \delta_{iN_q}\delta_{kN_q} -
\delta_{i1}\delta_{k1}$ reduces the symmetric part to pure boundary
terms (which cancel against the interface analysis below); the
antisymmetric part telescopes over the element if and only if the paired
contribution is a difference form $c_i - c_m$.

\begin{prop}[KEP volume flux]
With the vector-invariant momentum pairing (pointwise $\nabla K$ plus
central $K$ lifting), the antisymmetric part is a difference form --- and
$P_{\mathrm{adv}}$ receives no volume contribution --- if and only if the
two-point mass flux is the average of the nodal contravariant mass
fluxes,
\begin{equation}
M^\sharp\big(\boldsymbol{n}_a, \boldsymbol{n}_b;\, y_a, y_b\big)
= \tfrac12\big(\rho_a\, \uu_a\!\cdot\!\boldsymbol{n}_a
             + \rho_b\, \uu_b\!\cdot\!\boldsymbol{n}_b\big)
\;=\; \avg{\rho \tilde u}.
\end{equation}
The Kennedy--Gruber product form $\avg{\rho}\avg{\tilde u}$ is
\emph{not} kinetic-energy-compatible with this pairing.
\end{prop}

\subsection{Interface conditions}

Let $G = \rho\,\uu\cdot\hat n$ denote the nodal normal mass flux on a
face. Collecting the boundary terms of the volume analysis with the SAT
updates \eqref{eq:sat} and the $K$-lifting \eqref{eq:lift}, the
per-face-node production reduces to the compact form
\begin{equation}
P_{\mathrm{face}}
= s\!W\!J\;\jump{K}\,\big(\avg{G} - F^*_\rho\big)
\;+\; P_{\mathrm{pen}},
\label{eq:faceprod}
\end{equation}
where $P_{\mathrm{pen}}$ is the work of the velocity jump penalties.
Equation \eqref{eq:faceprod} dictates the interface set:

\begin{prop}[KEP interface set]
The interface production vanishes ($P_{\mathrm{face}} =
P_{\mathrm{pen}} \le 0$) with:
\begin{enumerate}
\item central interface mass flux $F^*_\rho = \avg{\rho\,\uu\cdot\hat n}
 = \avg{G}$, with \emph{no density penalty}: any added
 $-\alpha\jump{\rho}$ term inserts
 $+\,\alpha\,s\!W\!J\,\jump{\rho}\jump{K}$ into \eqref{eq:faceprod},
 which has no sign;
\item energy flux riding the same mass flux,
 $F^*_{\rho e} = \avg{e}\,F^*_\rho + \avg{p}\avg{\tilde u}$
 (the pressure-work part is the physical exchange channel and is
 excluded from $P_{\mathrm{adv}}$ by construction);
\item central traces for $K$ and $p$ in all lifting terms;
\item $\rho$-weighted velocity jump penalties applied to the owned
 covariant components,
 $g_{\mathrm{pen}}^\pm = \frac{\bar\lambda}{2}
 \frac{\avg{\rho}}{\rho^\pm}\,( u^\mp - u^\pm)$,
 whose total production is exactly
 $P_{\mathrm{pen}} = -\tfrac{\bar\lambda}{2}\, s\!W\!J\,
 \avg{\rho}\,\big|\jump{\uu}\big|^2 \;\le\; 0$.
\end{enumerate}
The Lamb term $\om\times\uu$ is kinetic-energy-inert pointwise for any
discrete $\om$. The resulting budget identity is metric-independent: it
closes identically on flat and hypsography-warped grids, with no discrete
geometric conservation law required.
\end{prop}

Numerically (operator-level verification on deterministic discontinuous
states, flat and Hughes--Jablonowski warped grids):
$|P_{\mathrm{adv}}| \lesssim 3\times10^{-15}\,P_{\mathrm{ref}}$ for the
KEP set, while the Kennedy--Gruber set violates the identity by 9--10
orders of magnitude, growing a further $10^3\times$ over terrain.

\section{Entropy-dissipative interfaces (\texttt{face\_set = :es})}
\label{sec:es}

Exact KEP does not control the thermodynamic (pressure-work) channel.
Introduce the mathematical entropy
$S = -\rho s/(\gamma - 1)$, $s = \ln(p\,\rho^{-\gamma})$. Holding $\rho$
(and the diagnostic $K$, $\Phi$) fixed, the entropy variable conjugate to
$\rho e$ follows from \eqref{eq:eos} by the chain rule:
\begin{equation}
v \;\equiv\; \frac{\partial S}{\partial(\rho e)}\bigg|_{\rho}
= -\frac{\rho}{\gamma-1}\,\frac{\partial s}{\partial (\rho e)}
= -\frac{\rho}{\gamma-1}\cdot\frac{\gamma-1}{p}
= -\frac{\rho}{p}.
\end{equation}
The \texttt{:es} interface set keeps every KEP central flux --- in
particular the mass flux stays central, so the kinetic-energy budget of
Section~\ref{sec:kep} is untouched --- and replaces the
entropy-indefinite Rusanov term on $\rho e$ by dissipation in $v$:
\begin{equation}
F^*_{\rho e} = F^{\mathrm{central}}_{\rho e}
 \;-\; \frac{\lambda}{2}\,\bar w\,\big(v^+ - v^-\big),
\qquad
\bar w = \frac{\bar p^{\,2}}{(\gamma - 1)\,\bar\rho} \;>\; 0,
\label{eq:esflux}
\end{equation}
with $\lambda$ the local maximum wave speed and overbars denoting face
arithmetic means.

\begin{prop}[Single-slot entropy dissipation]
Because \eqref{eq:esflux} modifies only the $\rho e$ slot and the SAT
update \eqref{eq:sat} is antisymmetric, the total entropy production of
the dissipative part is, exactly,
\begin{equation}
\Delta \dot S
= s\!W\!J\,\big(v^- - v^+\big)\cdot
  \Big(-\tfrac{\lambda}{2}\bar w\big(v^+ - v^-\big)\Big)\cdot(-1)
= -\frac{\lambda}{2}\,s\!W\!J\;\bar w\,\jump{v}^2 \;\le\; 0 .
\end{equation}
No matrix dissipation theory and no coupling to $K$ or $\Phi$ is
required. Near constant states, $\jump{v} \approx \bar\rho\,\jump{p}/
\bar p^{\,2}$, so \eqref{eq:esflux} reduces to Rusanov acting on the
internal-energy part of $\jump{\rho e}$; and the term is
kinetic-energy-inert since $K$ is diagnostic.
\end{prop}

\begin{rem}
This provides interface entropy \emph{dissipation}, not a full
entropy-stability proof: the KEP central fluxes produce entropy at
$O(\jump{\cdot}^3)$, and an entropy-\emph{conservative} volume flux
(Chandrashekar/Ranocha log-mean type) would break the exact VI-KEP
volume budget identity --- the two exactness properties compete at the
volume level. Empirically the modification is decisive: the
double-mountain baroclinic wave went from sub-hour top-of-atmosphere
pressure collapse to a full stable day with $\min p$ constant.
\end{rem}

\section{Well-balanced Exner pressure gradient}
\label{sec:exner}

The horizontal pressure-gradient force is discretized exclusively in a
reference-subtracted split Exner form (the plain $-\nabla_h p/\rho$ form
was removed after measuring it strictly worse in every configuration).
With Exner pressure $\Pi = (p/p_0)^{\Rd/\cp}$ and virtual potential
temperature $\theta = p/(\rho \Rd \Pi)$, the identity
$\dd p/\rho = \cp\,\theta\,\dd\Pi$ converts the PGF to
$-\cp\theta\nabla_h\Pi$. Choose a reference profile as a function of
pressure alone,
\begin{equation}
T_r(\Pi) = T_{\min} + \big(T_{\mathrm{sfc}} - T_{\min}\big)\,\Pi^{s},
\qquad
\theta_r = \frac{T_r}{\Pi},
\qquad
\Phi_r(\Pi) \;\text{with}\;
\frac{\dd \Phi_r}{\dd \Pi} = -\cp\,\theta_r,
\label{eq:ref}
\end{equation}
($s = 7$; $\Phi_r$ closed-form). By construction the reference is in
exact hydrostatic balance, $\cp\theta_r\nabla\Pi + \nabla\Phi_r = 0$,
\emph{pointwise as an algebraic identity in $p$} --- not as a competition
between two large discrete gradients. Splitting
$\theta = \theta_r + \theta'$ and symmetrizing the perturbation term
yields the discretized form
\begin{equation}
-\frac{1}{\rho}\nabla_h p - \nabla_h(K+\Phi)
\;\longrightarrow\;
-\nabla_h\big(K + \Phi - \Phi_r\big)
\;-\;\frac{\cp}{2}\Big[\theta'\,\nabla_h\Pi
 + \nabla_h\big(\theta'\Pi\big) - \Pi\,\nabla_h\theta'\Big],
\label{eq:exner}
\end{equation}
each strong gradient completed by a central lifting \eqref{eq:lift}. The
discrete terrain residual then scales with the $\theta$-perturbation
rather than with the full hydrostatic pressure variation along warped
coordinate surfaces.

\begin{rem}[Problem-matched reference]
The construction only helps if $T_r$ is close to the actual state. For
isothermal states $T \equiv T_0$, the matching choice is $T_r \equiv
T_0$, for which \eqref{eq:ref} integrates to $\Phi_r = -\cp T_0 \ln\Pi =
gz$ exactly and $\theta' \equiv 0$. The ClimaAtmos default $(T_{\min},
T_{\mathrm{sfc}}) = (220, 290)\,$K is $40\,$K off such states and voids
the cancellation entirely (measured: no reduction relative to the direct
form). The implementation therefore dispatches the reference per problem.
On the Hughes--Jablonowski double mountain, the matched form reduces the
$t=0$ momentum residual $11\times$ and removes a spurious
$5\,\mathrm{m\,s^{-1}\,h^{-1}}$ linear spin-up of the velocity field.
\end{rem}

\section{Discrete hydrostatic initialization}
\label{sec:hydro}

The staggered vertical balance seen by the $w$ equation
\eqref{eq:implicit} is the centered face relation
\begin{equation}
\frac{p_{v+1} - p_v}{\Delta z}
= -\,\frac{\rho_v + \rho_{v+1}}{2}\;
  \frac{\Delta(\Phi + K)}{\Delta z},
\label{eq:facebal}
\end{equation}
per interval $v \to v{+}1$ of each column. Analytic hydrostatic profiles
satisfy \eqref{eq:facebal} only to $O(\Delta z^2)$, which projects
$O(10\,\mathrm{m\,s^{-1}})$ spurious gravity-mode transients at
initialization if left uncorrected. Two exact discrete constructions are
used.

\paragraph{General profiles (sphere).}
Keeping $p$ analytic and solving \eqref{eq:facebal} for density gives the
recursion $\rho_{v+1} = -\rho_v - 2\,\Delta p/(g \Delta z)$, whose
homogeneous solution is $(-1)^v$: it deposits a level-alternating
(checkerboard) $\delta\rho$ of relative size $O\big((\Delta z/2H)^2\big)$.
Over terrain the checkerboard amplitude varies along coordinate surfaces
and acts as a horizontal PGF error $\nabla_h p \cdot \delta(1/\rho)$ ---
measured to be the dominant, discretization-form-independent terrain
imbalance. A smooth alternative for arbitrary $T(z)$ is the product
recursion
\begin{equation}
p_{v+1} = p_v\,\frac{1 - \beta_v}{1 + \beta_{v+1}},
\qquad
\beta_v = \frac{g\,\Delta z}{2 \Rd T_v},
\end{equation}
whose homogeneous solution is constant in $v$ (eigenvalue $+1$, no
checkerboard); it is recommended as the default and is exact per
interval on stretched grids.

\paragraph{Isothermal profiles (mountain-wave slab).}
With $\rho = p/(\Rd T_0)$, relation \eqref{eq:facebal} admits the exact
per-interval geometric solution
\begin{equation}
p_{v+1} = p_v\;\frac{1 - \delta_v/2}{1 + \delta_v/2},
\qquad
\delta_v = \frac{\Delta(\Phi + K)}{\Rd T_0}.
\label{eq:geom}
\end{equation}
The \emph{effective geopotential} $\Phi + K$ in \eqref{eq:geom} is
essential over steep terrain: the kinematic surface $w$ enters $K$
through \eqref{eq:fullK} (e.g.\ $\pm 13\,\mathrm{m\,s^{-1}}$ at slope
$0.65$), and balancing against $\Phi$ alone leaves an $O(\Delta K)$
vertical impulse. With $\Phi + K$, the $t=0$ vertical residual is machine
roundoff ($5\times10^{-11}$) even at $33^\circ$ slopes.

\section{Dissipation and damping inventory}
\label{sec:diss}

\begin{center}
\begin{tabular}{@{}lll@{}}
\toprule
Mechanism & Form & Notes\\
\midrule
$\kappa_4$ biharmonic & SIPG/LDG $\nabla_h^4$ on perturbations & terrain-aware:
acts on $(h_{\mathrm{tot}}, u, v)$ minus base state;\\
 & $\kappa_4 = f\cdot \Delta h^3/[(2N{+}1)^2\Delta t]$ & fraction $f$ of the
explicit stability cap\\
Velocity penalties & $\frac{\bar\lambda}{2}\frac{\avg{\rho}}{\rho}\jump{u}$ &
exact kinetic-energy sink (part of the KEP set)\\
\texttt{:es} interface & $-\frac{\lambda}{2}\bar w\jump{v}$ on $\rho e$ &
entropy-dissipative, kinetic-energy-inert\\
Divergence damping & $\nu_\delta \nabla_h(\nabla_h\!\cdot\uh)$ &
scale-selective on acoustic/divergent modes\\
$\nu_{\mathrm{vert}}$ & $\partial_z(\nu(z)\partial_z \uh)$ &
sponge-profile-weighted momentum deposition aloft\\
Rayleigh sponge & $-\beta(z)\,w$ (implicit) & $w$-only; ClimaAtmos strength
$\alpha_w = 1\,\mathrm{s^{-1}}$ admissible\\
Vertical transport & Lin--van Leer limiter & monotone explicit part of
HEVI advection\\
\bottomrule
\end{tabular}
\end{center}

\section{Results}
\label{sec:results}

\subsection{Operator-level verification}

On deterministic discontinuous states (spectral gradients of $C^0$
fields), flat and Hughes--Jablonowski-warped grids alike, the discrete
kinetic-energy budget closes to $3\times10^{-15}$ relative for the
\texttt{:kep} and \texttt{:es} sets, with sign-definite penalty
production; the isolated \texttt{:es} dissipation yields entropy
production $P_S = -1.6\times10^{9} < 0$ with mass fluxes bit-identical to
central.

\subsection{Baroclinic wave over double-mountain orography}

The Ullrich et al.\ baroclinic wave over the Hughes--Jablonowski analytic
double mountain (cubed sphere, \texttt{helem} 8, $30\,$km lid,
\texttt{:es} faces, matched Exner PGF) completes a stable simulated day:
mass conserved to $0.0$ and total energy to $-4.9\times10^{-16}$
relative over 1440 steps, $\min p$ constant throughout, with orographic
wave trains developing at $z = 2.5\,$km.

\subsection{Agnesi mountain waves}

The quasi-two-dimensional slab configuration
(Section~\ref{sec:terrain}; isothermal $T_0 = 250\,$K, $U_0 =
20\,\mathrm{m\,s^{-1}}$, ridge half-width $a = 25\,$km, 600~km periodic
domain, $2$~km elements, 750~m levels, implicit sponge above 15~km)
provides quantitative validation against linear mountain-wave theory.

\begin{figure}[htbp]
\centering
\includegraphics[width=0.85\textwidth]{mw_agnesi_linear_xz.png}
\caption{Steady vertical velocity $w(x,z)$ at $t = 4\,$h for the linear
Agnesi case, $h_0 = 25\,$m. Upstream-tilted phase lines (upward energy
propagation) with vertical wavelength $\lambda_z = 2\pi U_0/N =
6.4\,$km and amplitude at the linear scale $U_0 h_0/a =
0.02\,\mathrm{m\,s^{-1}}$ with the expected $\sqrt{\rho_0/\rho}$
amplification; two-dimensionality is preserved to $10^{-10}$.}
\label{fig:agnesi-linear}
\end{figure}

\begin{figure}[htbp]
\centering
\includegraphics[width=0.85\textwidth]{mw_agnesi_h250m_xz.png}
\caption{As Fig.~\ref{fig:agnesi-linear} but with $h_0 = 250\,$m
($10\times$ forcing) on the identical grid with no parameter retuning.
The wave structure is unchanged and the amplitude scales exactly
linearly ($0.34 \approx 10\times0.033\,\mathrm{m\,s^{-1}}$ below
15~km); the solution remains laminar ($N h_0/U_0 = 0.125$), whereas with
a $300\times$ weaker explicit sponge the same case breaks aloft through
$\sqrt{\rho_0/\rho}$ amplification and is fully turbulent by
$t \approx 3\,$h.}
\label{fig:agnesi-h250}
\end{figure}

At the steep limit $h_0 = a = 2\,$km ($N h_0/U_0 \approx 2$, $33^\circ$
slopes) the $\Phi+K$-balanced initialization of
Section~\ref{sec:hydro} holds the $t=0$ vertical residual at roundoff;
the remaining obstacle is the impulsive uniform-wind start itself
($\approx 4\,\mathrm{m\,s^{-2}}$ horizontal adjustment shock), which
requires a ramped wind-forcing initialization rather than further
discretization work.

\end{document}
